\documentclass[conference]{IEEEtran}
\IEEEoverridecommandlockouts

\usepackage{cite}
\usepackage{amsmath,amssymb,amsfonts}
\usepackage{graphicx}
\usepackage{textcomp}
\usepackage{xcolor}
\usepackage{booktabs}
\usepackage{array}
\usepackage{url}
\usepackage{hyperref}
\usepackage{multirow}

\hypersetup{
    colorlinks=true,
    linkcolor=black,
    citecolor=black,
    urlcolor=black
}

\def\BibTeX{{\rm B\kern-.05em{\sc i\kern-.025em b}\kern-.08em
    T\kern-.1667em\lower.7ex\hbox{E}\kern-.125emX}}

\begin{document}

\title{Beyond the Perimeter: A Survey of Enterprise VPN Technologies\\
and a Proposed Quantum-Resilient Adaptive Zero Trust Architecture}

\author{
    \IEEEauthorblockN{[Author 1 Name]}
    \IEEEauthorblockA{\textit{Department of Computer Science} \\
    \textit{Habib University}\\
    Karachi, Pakistan \\
    student1@stu.habib.edu.pk}
    \and
    \IEEEauthorblockN{[Author 2 Name]}
    \IEEEauthorblockA{\textit{Department of Computer Science} \\
    \textit{Habib University}\\
    Karachi, Pakistan \\
    student2@stu.habib.edu.pk}
    \and
    \IEEEauthorblockN{[Author 3 Name]}
    \IEEEauthorblockA{\textit{Department of Computer Science} \\
    \textit{Habib University}\\
    Karachi, Pakistan \\
    student3@stu.habib.edu.pk}
}

\maketitle

\begin{abstract}
Enterprise VPN security is at a turning point. The technologies
that have underpinned corporate network access for the past two
decades --- GRE tunneling, IPSec encryption, and MPLS switching
--- were not designed for a world of cloud-native workloads,
ubiquitous remote work, and quantum-capable adversaries. This
paper does two things. First, we survey six enterprise
connectivity paradigms (GRE, IPSec, WireGuard, MPLS L3VPN,
SD-WAN, and ZTNA) with a focus on the threat dimensions most
relevant today: lateral movement risk, quantum vulnerability,
and performance overhead. Rather than treating these as
equivalent considerations, we argue that quantum key exposure
and post-authentication lateral movement now represent the two
defining security failures of the traditional VPN model --- a
claim we support through a structured threat-surface taxonomy
and performance data drawn from recent peer-reviewed studies.
Second, we propose a conceptual integration framework, the
\textit{Quantum-Resilient Adaptive Zero Trust Architecture}
(QRA-ZTA), that layers post-quantum cryptographic key exchange
(NIST FIPS-203 ML-KEM), eBPF-accelerated policy enforcement,
zero trust continuous verification (per NIST SP~800-207), and
AI-driven traffic anomaly detection into a coherent architecture
that can be deployed atop existing enterprise transport
infrastructure. The goal is not to replace what organizations
have built but to give those investments the security properties
they were never designed to have.
\end{abstract}

\begin{IEEEkeywords}
VPN, GRE, IPSec, WireGuard, MPLS, SD-WAN, ZTNA, zero trust,
post-quantum cryptography, ML-KEM, eBPF, encrypted traffic
anomaly detection, HNDL attacks, enterprise network security,
QRA-ZTA
\end{IEEEkeywords}

% ---------------------------------------------------------------
\section{Introduction}

There is a particular irony in how enterprise VPN security
actually works in practice. The whole effort --- the cryptographic
handshake, the key exchange, the mutual authentication --- goes
into getting a device through the tunnel. Once it is through,
that device is on the network. It can reach servers, databases,
shared drives. The implicit assumption is that anything inside
the tunnel belongs there. This is the castle-and-moat model, and
it has been the foundational logic of enterprise network security
for thirty years.

The problem is that attackers understand this model better than
most defenders do. Compromising a single VPN credential, whether
through phishing, password spraying, or exploiting an unpatched
gateway vulnerability, grants the kind of broad internal
reachability that turns a single foothold into a domain-wide
breach. The Cybersecurity Insiders 2024 VPN Risk Report found
that 56\% of enterprises had experienced at least one cyberattack
in the preceding year that exploited a VPN vulnerability
directly~\cite{cybersec2024}. Ransomware operators have made
VPN gateways their preferred initial access vector precisely
because of this dynamic.

A second threat is accumulating more quietly. State-sponsored
actors are conducting what the intelligence community calls
``Harvest Now, Decrypt Later'' (HNDL) attacks: intercepting and
storing encrypted VPN traffic today, with the intention of
decrypting it once quantum computers become capable enough.
The key exchange algorithms that secure every mainstream VPN
--- Diffie-Hellman in IPSec, Curve25519 in WireGuard, ECDH in
TLS-based solutions --- are all solvable in polynomial time by
Shor's algorithm on a sufficiently large quantum machine. The
timeline for when such machines might exist is genuinely
uncertain, but the intelligence community estimates the window
at ten to fifteen years~\cite{nist_fips203}. The traffic being
collected right now will still be sensitive then. This is not a
hypothetical future problem; the collection is happening today.

Against this backdrop, this paper argues that the incremental
hardening of legacy VPN stacks is no longer sufficient as a
security strategy. What is needed instead is a clear-eyed
assessment of what each VPN technology actually provides and
where it structurally fails, followed by a composable
architecture that addresses those failures without requiring
enterprises to abandon the infrastructure they have built.

We organize this argument in four parts. Section~\ref{sec:background}
establishes a taxonomy of the six paradigms we examine. The
survey in Section~\ref{sec:relatedwork} covers each technology
in enough depth to expose its specific threat surface.
Section~\ref{sec:threatscape} consolidates that analysis into
a structured gap assessment, including quantitative performance
data. Section~\ref{sec:framework} then proposes the QRA-ZTA
framework as a response to those gaps.

% ---------------------------------------------------------------
\section{Organizing the Landscape: A Three-Dimension Taxonomy}
\label{sec:background}

Before comparing technologies, it helps to agree on what we are
comparing them on. Prior surveys of enterprise VPN have typically
organized technologies by operational layer (L2 vs.\ L3) or
management model (customer-managed vs.\ provider-managed). We
retain both of those dimensions but add a third: trust model.
This distinguishes between perimeter-based architectures, where
trust is established once at connection time and implicitly
extended throughout the session, and zero-trust architectures,
where every access request is evaluated independently regardless
of how the user got onto the network~\cite{nist800207}.

\begin{table}[htbp]
\caption{Taxonomy of Enterprise Connectivity Paradigms}
\label{tab:taxonomy}
\centering
\renewcommand{\arraystretch}{1.2}
\begin{tabular}{lccc}
\toprule
\textbf{Technology} & \textbf{Management} & \textbf{Layer} & \textbf{Trust Model} \\
\midrule
GRE         & Customer  & L3      & Perimeter \\
IPSec       & Customer  & L3      & Perimeter \\
WireGuard   & Customer  & L3      & Perimeter \\
MPLS L3VPN  & Provider  & L3      & Perimeter \\
SD-WAN      & Hybrid    & L3      & Perimeter/Hybrid \\
ZTNA        & Cloud     & App     & Zero Trust \\
\bottomrule
\end{tabular}
\end{table}

The trust model dimension is why this taxonomy matters more than
a simple feature comparison. The first five technologies in
Table~\ref{tab:taxonomy} all share perimeter-based trust, which
means they share the same structural vulnerability to lateral
movement after credential compromise. ZTNA breaks from this
model but, as we will argue, introduces its own gaps --- notably
the absence of quantum-resistant key exchange and the limitation
of not providing site-to-site transport. The QRA-ZTA framework
in Section~\ref{sec:framework} is designed to fill the
intersection of those gaps.

% ---------------------------------------------------------------
\section{Survey of Enterprise VPN Technologies}
\label{sec:relatedwork}

\subsection{GRE: Flexible but Fundamentally Insecure}

GRE, standardized in RFC~2784~\cite{rfc2784}, encapsulates one
network layer protocol inside another. In practice this means
you can carry IPv6 over an IPv4 backbone, run multicast across
a unicast WAN, or maintain OSPF adjacencies between routers
that have no direct Layer-2 connection. Fahriza et
al.~\cite{fahriza2026} document a three-branch enterprise
deployment using GRE for IPv6-over-IPv4 tunneling with OSPF
routing, and Huda and Andriyani~\cite{huda2025} examine GRE's
integration with OSPF multi-area configurations. Both confirm
GRE's practical value as a routing and encapsulation primitive.

What GRE does not do is provide any security. There is no
encryption, no authentication, and no replay protection. Uddin
et al.~\cite{uddin2021} are direct about this: GRE alone is
unsuitable for any environment where data confidentiality
matters. Ogudo~\cite{ogudo2019} reaches the same conclusion,
noting that GRE over a public network without an overlaid
security protocol creates a completely open interception surface.
GRE belongs in the toolkit as a routing mechanism, not a
security one.

\subsection{IPSec: The Enterprise Standard, and Its Problems}

IPSec (RFC~4301~\cite{rfc4301}) is the closest thing enterprise
networking has to a universal standard for encrypted tunnels.
It provides confidentiality through ESP, integrity through keyed
hashes, and mutual authentication through IKE. Deployed in
tunnel mode with ESP, it encrypts the entire original IP packet
inside a new outer envelope. Nair~\cite{nair2024} identifies
IKEv2 with AES-256-GCM as the correct configuration for new
deployments, and Rosu et al.~\cite{rosu2012} validate IPSec at
enterprise scale in a large Romanian corporate network.

The operational vulnerabilities are real and well-studied.
Sankaran~\cite{sankaran2023} documents IKEv1 downgrade attacks,
exploitable cipher negotiations, and SA exhaustion as active
threats against deployed IPSec configurations. IKEv1 should be
considered deprecated at this point. More fundamentally, IPSec's
key exchange relies on the discrete logarithm problem, which
Shor's algorithm solves efficiently on a quantum
machine~\cite{nist_fips203}. This is not an academic concern.
Every enterprise running IPSec today is generating traffic that
HNDL attackers may already be collecting.

There is also the operational complexity problem. The IKE
negotiation stack is large, the configuration options are
numerous, and misconfiguration is common. The connection model
requires active state management, which means mobile users who
roam between networks frequently drop and re-establish tunnels.
These are not disqualifying limitations, but they are real costs.

\subsection{WireGuard: A Better Mousetrap}

WireGuard~\cite{donenfeld2017} was designed as an explicit
rejection of IPSec's complexity. Where IPSec negotiates
cryptographic algorithms, WireGuard picks one suite and uses
it: Curve25519 for key agreement, ChaCha20-Poly1305 for
authenticated encryption, and BLAKE2s for hashing. There is
nothing to negotiate and therefore nothing to downgrade.
Dowling and Paterson~\cite{dowling2018} formally verified the
WireGuard handshake using the eCK security model and confirmed
it achieves forward secrecy and resistance to key compromise
impersonation --- properties that IPSec configurations often
get wrong through misconfiguration.

The codebase is roughly 4,000 lines~\cite{calmops2026}, a
fraction of OpenVPN's 600,000-plus lines. A smaller codebase
is not just an aesthetic preference; it is a measurable reduction
in attack surface and a practical prerequisite for formal
verification. Performance reflects the design: independent
benchmarks show WireGuard at 92--95\% of native bare-metal
bandwidth, compared to 85--90\% for IKEv2/IPSec in comparable
software deployments, with latency overhead of 1--3~ms versus
10--15~ms~\cite{zhuque2026,zenarmor2023}. Tanveer and
Ashfaq~\cite{tanveer2023} identify it as the strongest
candidate for modern software-based VPN from a combined security
and performance standpoint.

The qualification is important: WireGuard's Curve25519 key
exchange is just as quantum-vulnerable as IPSec's DH groups.
This is WireGuard's most significant unresolved problem for
any organization thinking beyond the next five years.

\subsection{GRE over IPSec: A Necessary Compromise}

Running GRE inside an IPSec tunnel is the standard workaround
for a real limitation: IPSec cannot carry multicast traffic
or routing protocol updates natively, but many enterprise WAN
designs depend on OSPF or EIGRP adjacencies. GRE-over-IPSec
combines GRE's routing flexibility with IPSec's encryption.
Fahriza et al.~\cite{fahriza2026} evaluate this in a three-branch
deployment and find it operationally sound with reasonable QoS
outcomes. Kumar et al.~\cite{kumar2020} specifically examine
multicast transport through GRE before IPSec encryption.

The cost is header overhead: at minimum 4 bytes of GRE plus
28 bytes of ESP overhead per packet, which reduces effective
MTU and adds CPU load. For throughput-sensitive paths this
matters and needs to be factored into capacity planning.
The quantum vulnerability of the underlying IPSec key exchange
carries over completely.

\subsection{MPLS L3VPN: Carrier-Grade Isolation}

MPLS (RFC~3031~\cite{rfc3031}) routes packets based on
short labels rather than full IP lookups, enabling the
provider to build logically isolated VPNs over shared physical
infrastructure. The L3VPN variant (RFC~4364~\cite{rfc4364})
uses per-customer VRF tables on Provider Edge routers and
MP-BGP to distribute routing information between them. Route
Distinguishers make overlapping customer address spaces
globally unique, and Route Targets control which routes cross
VPN boundaries. Mungse~\cite{mungse2023} gives a detailed
treatment of VRF operation. Fayed et al.~\cite{fayed2024}
validate three progressively complex L3VPN models in GNS3,
including an overlapping VPN and a central services
architecture where a shared server farm is reachable from
multiple customer VRFs through controlled route leaking.

MPLS has genuine strengths: native QoS through EXP bits in
the label header, traffic engineering via RSVP-TE, and a
core that stays clean regardless of how many customers
you add. The central services model Fayed et al.\ demonstrate
is difficult to replicate with GRE or IPSec at comparable
scale.

However, MPLS does not encrypt anything. The isolation is
logical, not cryptographic. Traffic traversing shared provider
infrastructure is readable by anyone with access to that
infrastructure. For organizations with regulatory requirements
for data-in-transit encryption, MPLS must be combined with
an IPSec overlay --- which then inherits IPSec's quantum
vulnerability. The dependence on provider infrastructure also
means MPLS cannot extend natively into public cloud
environments.

\subsection{SD-WAN: Operational Flexibility, Evolving Security}

SD-WAN decouples WAN control from the underlying transport,
allowing a central controller to aggregate multiple link types
--- MPLS circuits, broadband Internet, LTE --- and route traffic
dynamically based on application requirements and real-time
link quality~\cite{tahenni2024}. The market has grown rapidly
for a straightforward reason: it substantially reduces WAN costs
compared to all-MPLS designs while providing acceptable
performance for most traffic, projected at a CAGR of 10.72\%
through 2031~\cite{fortinet_sdwan}.

Tahenni and Merazka~\cite{tahenni2024} find that MPLS still
wins on deterministic latency for mission-critical real-time
traffic, which is why the dominant enterprise pattern is
hybrid: SD-WAN augmenting MPLS rather than replacing it
outright. SD-WAN also resolves the $O(N^2)$ tunnel scaling
problem of manually configured IPSec: the controller provisions
tunnels automatically, and WireGuard is increasingly adopted
as the encryption underlay. The security concern is the
controller itself, which concentrates the attack surface for
the entire WAN overlay into a single logical component.

\subsection{ZTNA: The Right Idea, Incomplete Alone}

ZTNA applies the zero trust model~\cite{nist800207} to access
control: no implicit trust after authentication, no broad
network reachability, access granted only to the specific
application the user is requesting right now. The NIST
SP~800-207 framework organizes this around three components ---
Policy Engine (PE), Policy Administrator (PA), and Policy
Enforcement Point (PEP) --- that together ensure a resource
remains invisible to anyone who has not been explicitly granted
access to it in the current session.

The case for ZTNA is empirically strong. Gartner projected
that 70\% of new remote-access deployments by 2025 would be
ZTNA-based~\cite{gartner2023}, and the 56\% enterprise VPN
breach rate in 2024~\cite{cybersec2024} explains why organizations
are moving in that direction. Eliminating lateral movement
structurally --- by making it impossible for a compromised
credential to reach anything beyond its authorized application
--- is a more complete solution to the post-authentication
trust problem than any amount of network segmentation.

But ZTNA is not a complete answer. It does not provide
site-to-site WAN transport. Its key exchange is typically
classical and therefore quantum-vulnerable. And it shifts the
attack surface to the identity provider: compromise the IdP
and you may compromise everything. We do not raise these points
to dismiss ZTNA but to identify precisely where the gaps are
that a more complete architecture needs to fill.

% ---------------------------------------------------------------
\section{Where Things Actually Stand: Gap Analysis}
\label{sec:threatscape}

\subsection{Threat Surface Taxonomy}

Table~\ref{tab:threats} maps the primary threat vectors for
each technology across three categories: passive threats
(interception without modification), active threats (injection,
impersonation, disruption), and quantum threats (attacks that
become feasible with a cryptographically relevant quantum
computer).

\begin{table}[htbp]
\caption{Threat Surface Taxonomy by Technology}
\label{tab:threats}
\centering
\renewcommand{\arraystretch}{1.2}
\begin{tabular}{p{1.3cm}p{6.2cm}}
\toprule
\textbf{Technology} & \textbf{Primary Threat Vectors} \\
\midrule
GRE & Passive: full payload interception. Active: packet injection, endpoint spoofing. Quantum: N/A (no cryptography to attack). \\
\midrule
IPSec & Active: IKEv1 downgrade~\cite{sankaran2023}, cipher negotiation exploitation, SA exhaustion DoS, PSK theft. Quantum: HNDL against DH/ECDH key exchange. \\
\midrule
WireGuard & Smallest overall attack surface; no downgrade possible. Active: roaming endpoint impersonation (requires stolen private key). Quantum: HNDL against Curve25519. \\
\midrule
MPLS L3VPN & Active: label spoofing (requires provider infrastructure access), provider insider threat. Passive: unencrypted payload interception on shared provider infrastructure. \\
\midrule
SD-WAN & Active: controller plane compromise, edge device impersonation. Quantum: HNDL against IPSec/WireGuard underlay key exchanges. \\
\midrule
ZTNA & Active: IdP compromise, device posture bypass, OAuth token theft, cloud broker availability (DoS). Quantum: HNDL if PQC not integrated into handshake. \\
\bottomrule
\end{tabular}
\end{table}

\subsection{The Lateral Movement Problem}

This is worth stating plainly. In GRE, IPSec, WireGuard, and
MPLS deployments, a successfully authenticated user gets
IP-level reachability to a network segment, not to a specific
application. The tunnel is the access control. If the device
is later compromised, or if the credentials were stolen and
replayed from a different machine, the attacker has the same
access the legitimate user had. This is how ransomware operators
escalate from a phished credential to a domain controller.
It is not a configuration error; it is a consequence of how
these protocols were designed~\cite{cybersec2024,tanveer2023}.

ZTNA solves this through application-granular enforcement,
but only for user-to-application remote access. Site-to-site
connectivity, which still makes up a large fraction of
enterprise WAN traffic, does not have a ZTNA equivalent.

\subsection{The Quantum Horizon}

Two distinct quantum threats apply to VPN cryptography, and
they require different responses. Shor's algorithm solves the
discrete logarithm and integer factorization problems that
underlie RSA, Diffie-Hellman, and ECC in polynomial time,
rendering classical key exchange transparent~\cite{nist_fips203}.
Grover's algorithm provides a quadratic speedup for symmetric
key search, which effectively halves symmetric key security:
AES-128 drops to 64 effective bits under quantum attack,
making AES-256 the practical minimum going forward.

NIST finalized FIPS-203~\cite{nist_fips203} in August 2024,
standardizing the Module-Lattice Key Encapsulation Mechanism
(ML-KEM, formerly CRYSTALS-Kyber) for key encapsulation. ML-KEM
is grounded in the hardness of the Learning With Errors (LWE)
problem over module lattices, which is believed to remain
intractable for quantum computers. RFC~9370~\cite{rfc9370}
provides the IKEv2 protocol mechanism for running multiple key
exchanges within a single negotiation, giving IPSec deployments
a concrete migration path to hybrid PQC without abandoning
the existing protocol.

\subsection{Quantitative Performance Summary}

Table~\ref{tab:perf} synthesizes benchmark data from recent
empirical studies~\cite{zhuque2026,zenarmor2023,tahenni2024}.
These figures apply to software-based implementations on
general-purpose hardware without hardware crypto acceleration.
On machines with AES-NI, IPSec with AES-GCM can narrow the
throughput gap with WireGuard substantially.

\begin{table}[htbp]
\caption{Performance Benchmarks, Software Implementations}
\label{tab:perf}
\centering
\renewcommand{\arraystretch}{1.2}
\begin{tabular}{lccc}
\toprule
\textbf{Protocol} & \textbf{Throughput} & \textbf{Latency} & \textbf{CPU} \\
                  & (\% of native) & overhead (ms) & overhead \\
\midrule
WireGuard         & 92--95\%  & 1--3    & Low \\
IKEv2/IPSec       & 85--90\%  & 10--15  & Medium \\
OpenVPN           & 70--80\%  & 15--25  & High \\
GRE (unencrypted) & $\sim$99\%& $<$1    & Minimal \\
GRE over IPSec    & 83--88\%  & 11--16  & Medium \\
\bottomrule
\end{tabular}
\end{table}

% ---------------------------------------------------------------
\section{Proposed Framework: The QRA-ZTA}
\label{sec:framework}

The gap analysis in Section~\ref{sec:threatscape} leads to a
fairly specific architectural requirement. An enterprise access
architecture that addresses the current threat landscape needs
to do at least four things that no single existing technology
does: replace classical key exchange with quantum-resistant
equivalents, enforce application-granular access control with
continuous re-verification, process policy at kernel speed
without becoming a performance bottleneck, and detect threats
inside encrypted tunnels without decrypting them. The
Quantum-Resilient Adaptive Zero Trust Architecture (QRA-ZTA)
is our attempt to specify how those requirements can be met
together, in a way that organizations can layer atop their
existing infrastructure incrementally.

The framework has four strata. Fig.~\ref{fig:qrazta} shows
how they relate.

\begin{figure}[htbp]
\centering
\begin{tabular}{|c|}
\hline
\textbf{Stratum IV: Intelligence Layer} \\
AI Anomaly Detection on Encrypted Traffic Metadata \\
(eBPF Flow Statistics $\rightarrow$ Autoencoder/LSTM Engine) \\
\hline
\textbf{Stratum III: Identity and Trust Layer} \\
ZTNA Policy Engine / Policy Administrator / PEP \\
Continuous Verification, Least-Privilege Access \\
\hline
\textbf{Stratum II: Programmable Data Plane} \\
eBPF/XDP Kernel-Space Policy Enforcement \\
AES-256-GCM Symmetric Datapath (PQC-derived keys) \\
\hline
\textbf{Stratum I: Cryptographic Foundation} \\
Hybrid PQC Handshake: X25519 + ML-KEM (FIPS-203) \\
User-Space Key Establishment Daemon \\
\hline
\textbf{Transport Substrate (Existing Infrastructure)} \\
GRE / IPSec / WireGuard / MPLS / SD-WAN \\
\hline
\end{tabular}
\caption{The QRA-ZTA layer model. Each stratum addresses a
distinct threat class while interfacing with the transport
substrate organizations already operate.}
\label{fig:qrazta}
\end{figure}

\subsection{Stratum I: Hybrid Post-Quantum Key Exchange}

The quantum vulnerability of classical VPN key exchange cannot
be patched; it has to be replaced. Our approach is the hybrid
model standardized in RFC~9370~\cite{rfc9370}, which runs a
classical key exchange (X25519 or ECDH-P256) and an ML-KEM
encapsulation in the same handshake, deriving a shared secret
from both. The resulting session key is quantum-safe if either
component holds: if ML-KEM's lattice foundations are later
found to have an unexpected weakness, the classical component
still protects against classical adversaries. If a quantum
computer cracks X25519, the ML-KEM component makes the
captured ciphertext useless to an HNDL attacker.

This is not a hypothetical deployment model. Cloudflare has
run X25519+Kyber-768 hybrid in production TLS~1.3 across
its network~\cite{cloudflare_pqc}. Enterprise firewall vendors
have begun integrating ML-KEM into IKEv2 negotiation. The
mechanism exists and is being deployed; what we are proposing
is applying it systematically as the cryptographic foundation
of an integrated access architecture.

One practical constraint worth noting: ML-KEM public keys are
1,184 bytes for ML-KEM-768, compared to 32 bytes for Curve25519.
This increases handshake size meaningfully and can cause
fragmentation or timeout issues with middleboxes that
aggressively limit packet sizes. The computationally intensive
handshake is handled by a dedicated user-space daemon rather
than being run in kernel space, precisely because the instruction
count and complexity would exceed what in-kernel environments
can safely execute. Once the shared secret is established and
the symmetric AES-256-GCM session keys derived, those keys are
pushed down into the data-plane component described next.

\subsection{Stratum II: eBPF-Based Data Plane Enforcement}

Traditional Linux packet processing traverses the full TCP/IP
stack. For each packet this means context switches, sequential
iptables rule evaluation, and routing table lookups. At
enterprise traffic volumes, this overhead is measurable and
real. eBPF changes the model fundamentally: sandboxed programs
run directly within the kernel, attached to the eXpress Data
Path (XDP) hook at the NIC driver level, processing packets
before they reach the standard network stack~\cite{ebpf_networking}.

In the QRA-ZTA data plane, the session keys from Stratum I
are stored in eBPF Maps --- efficient in-kernel key-value
structures accessible by both user-space daemons and eBPF
programs. XDP-attached programs consult these maps to decrypt
incoming packets and enforce per-flow forwarding rules derived
from the access policies in Stratum III. The entire operation
bypasses iptables. The performance gains are significant:
eBPF-based networking has been shown to achieve packet
forwarding throughput comparable to dedicated hardware, with
latency reductions of an order of magnitude versus
iptables-based approaches~\cite{ebpf_networking}.

There is a security benefit beyond performance. The kernel
verifier statically analyzes every eBPF program before
execution, confirming it terminates, does not loop, and
maintains strict memory safety. No eBPF program can destabilize
the kernel host --- a property loadable kernel modules do not
share. This makes the data plane a safer place to enforce
security policy than a loadable module would be.

\subsection{Stratum III: Continuous Zero Trust Verification}

The zero trust model defined in NIST SP~800-207~\cite{nist800207}
is built on three components: the Policy Engine (PE), which
evaluates access requests; the Policy Administrator (PA), which
issues session tokens based on the PE's decisions; and the
Policy Enforcement Point (PEP), which gates data-plane access.
In QRA-ZTA, the PEP is the eBPF program in Stratum II. The PE
evaluates a composite trust score for each session request,
drawing on user identity (certificate or external IdP), device
posture (patch level, EDR agent status), contextual signals
(geography, time, access history), and threat intelligence.
The PA translates the PE's verdict into short-lived session
tokens and updates the eBPF Maps with per-flow forwarding
rules scoped to the specific application endpoint, not to a
network segment.

The critical difference from a conventional VPN is what happens
after the session is established. QRA-ZTA does not grant
ongoing implicit trust. Tokens expire on a schedule, requiring
periodic re-evaluation by the PE. More usefully, the anomaly
detection output from Stratum IV feeds directly into the PE:
a behavioral anomaly can trigger an immediate trust score
downgrade, prompting re-authentication or, for high-confidence
threat signals, immediate session revocation. The enforcement
is at the eBPF level, so revocation happens in the kernel fast
path with no application-layer round trip.

This design eliminates lateral movement structurally. A
compromised credential in a QRA-ZTA deployment grants access
to the specific application that credential was authorized
to reach, in the current session, not to the network segment
that application lives in. There is no path from a single
credential compromise to a domain-wide breach without
separately compromising additional credentials and passing
continuous behavioral analysis.

\subsection{Stratum IV: AI Anomaly Detection on Encrypted Traffic}

End-to-end encryption is a security benefit with an operational
cost: it blinds the traditional intrusion detection tools that
depend on reading packet payloads. Deep packet inspection does
not work on properly encrypted traffic, and attempting to
break encryption at a middlebox for inspection purposes
undermines the end-to-end security model and creates new
exposure points~\cite{ai_anomaly_det}. This is a genuine
tension, and it is one that adversaries have learned to
exploit --- encrypted C2 channels and exfiltration over HTTPS
are routine precisely because they evade payload-based
detection.

The response in QRA-ZTA is to shift from payload analysis to
behavioral analysis. Encrypted traffic has observable
characteristics that do not require reading the payload:
packet lengths, inter-arrival timing, flow duration, byte
volume distributions, and burst patterns all differ in
meaningful ways across different traffic types. A ransomware
beacon, a file exfiltration stream, and a video call have
very different behavioral profiles even when all three are
encrypted with AES-256.

eBPF programs in the data plane continuously compute these
statistics in kernel space, without copying payload data to
user space. The low overhead of this approach is the key
enabler: traditional packet capture tools like libpcap introduce
unacceptable performance penalties at enterprise traffic
volumes, but eBPF runs in the kernel fast path with negligible
impact. Aggregated statistics are pushed to a user-space AI
engine via ring buffers. The engine runs an ensemble: supervised
classifiers (Random Forest or XGBoost) handle high-confidence
classification of known threat patterns, while unsupervised
autoencoders and LSTM networks model per-user, per-application
behavioral baselines and flag deviations that do not match
any known signature~\cite{ai_anomaly_det}.

The DeSFAM framework~\cite{desfam}, which correlates eBPF
network telemetry with host-level system call traces, shows
what this approach can achieve: 96\% detection accuracy and
an F1-score of 0.92 for container-level threats, without
payload inspection. When QRA-ZTA's intelligence layer flags
an anomaly, the signal goes to the Stratum III Policy Engine,
which can downgrade the session trust score and trigger
re-verification. The loop is closed: telemetry drives
policy, policy drives enforcement, enforcement is in the
kernel fast path.

\subsection{How QRA-ZTA Layers onto Existing Infrastructure}

We want to be direct about what this framework is and what it
is not. It is not a proposal to rip out existing GRE, IPSec,
MPLS, or SD-WAN infrastructure and replace it. That is
operationally unrealistic for most enterprises, and it is
not necessary to realize the security benefits the framework
provides.

The framework is designed to be applied incrementally. An
organization running IKEv2/IPSec today can upgrade to hybrid
PQC key exchange using RFC~9370~\cite{rfc9370} without changing
anything else, addressing the HNDL threat against recorded
traffic immediately. The eBPF data plane can be introduced
alongside existing iptables configurations as an acceleration
layer for high-priority traffic. ZTNA policy enforcement in
Stratum III can be applied to the highest-value applications
first --- financial systems, HR databases, administrative
interfaces --- while GRE/IPSec/MPLS continues to handle
lower-priority traffic.

For MPLS deployments where payload encryption is absent,
Stratum I hybrid PQC key exchange applied as an IPSec overlay
addresses both the plaintext exposure and the quantum threat
in a single step. SD-WAN deployments benefit from Stratum I
applied to the existing WireGuard or IPSec underlay, with
Stratum IV anomaly detection applied to the controller's
telemetry stream. WireGuard deployments are the most
straightforward: the symmetric key infrastructure is already
in place, and only the key exchange requires upgrading once
a standardized WireGuard PQC extension becomes
available~\cite{hu2020pq}.

% ---------------------------------------------------------------
\section{Cross-Technology Comparative Analysis}
\label{sec:discussion}

Table~\ref{tab:comparison} extends the comparison across the
six technologies with the dimensions that the gap analysis in
Section~\ref{sec:threatscape} identified as most consequential.
The lateral movement and quantum resistance rows are not found
in prior comparative surveys of enterprise VPN; we add them
here because they are the two threat dimensions that most
clearly differentiate the current security requirements from
those of a decade ago.

\begin{table*}[htbp]
\caption{Extended Comparative Summary Including Security-Critical Dimensions}
\label{tab:comparison}
\centering
\renewcommand{\arraystretch}{1.3}
\begin{tabular}{lcccccc}
\toprule
\textbf{Criterion} & \textbf{GRE} & \textbf{IPSec} & \textbf{WireGuard}
  & \textbf{MPLS L3VPN} & \textbf{SD-WAN} & \textbf{ZTNA} \\
\midrule
Encryption          & None      & Strong    & Strong      & None$^{a}$    & Configurable & Strong \\
Initial auth.       & None      & Strong    & Strong      & Implicit      & Strong       & Strong \\
Continuous auth.    & No        & No        & No          & No            & No           & Yes \\
Lateral movement    & High      & High      & High        & High          & High         & Eliminated \\
Quantum-resistant   & No        & No$^{b}$  & No$^{b}$    & N/A           & No$^{b}$     & No$^{b}$ \\
Routing support     & Full      & Limited   & Limited     & Full (BGP)    & Full         & N/A \\
Native QoS          & No        & No        & No          & Yes           & Dynamic      & N/A \\
Scalability         & Low       & Low       & Medium      & High          & High         & High \\
Config. complexity  & Low       & Medium    & Low         & High          & Low          & Medium \\
Cloud readiness     & Medium    & Medium    & High        & Low           & Native       & Native \\
Throughput          & $\sim$99\%& 85--90\%  & 92--95\%    & N/A           & Variable     & N/A \\
Latency overhead    & $<$1~ms   & 10--15~ms & 1--3~ms     & Deterministic & Dynamic      & Per-session \\
\bottomrule
\multicolumn{7}{l}{\small $^{a}$MPLS provides logical isolation, not cryptographic encryption.} \\
\multicolumn{7}{l}{\small $^{b}$Hybrid PQC deployable via RFC~9370~\cite{rfc9370} for IPSec, WireGuard, SD-WAN underlays.} \\
\end{tabular}
\end{table*}

Reading across the lateral movement and quantum resistance rows
makes the architectural gap visible. No existing technology
handles both. ZTNA eliminates lateral movement but is
quantum-vulnerable. IPSec and WireGuard can be made
quantum-resistant through hybrid key exchange but do nothing
about post-authentication lateral movement. This is precisely
the gap QRA-ZTA addresses by combining Stratum I PQC key
exchange with Stratum III zero trust enforcement.

\subsection{Practical Selection Guidance}

For organizations that need to make technology decisions now,
the analysis supports the following guidance.

Security-first organizations --- financial institutions, healthcare,
defense contractors --- should apply RFC~9370 hybrid PQC key
exchange to existing IPSec site-to-site tunnels immediately.
The HNDL threat is active today, and recorded traffic cannot
be retroactively protected after a quantum computer becomes
available. For user-to-application access, ZTNA should replace
remote-access VPN for all privileged and sensitive resource access.

Performance-critical organizations should retain MPLS for
real-time, latency-sensitive traffic where its deterministic
QoS properties are genuinely required. WireGuard is the right
choice for all other site-to-site connectivity, offering the
best software-implementation performance by a meaningful margin.

Cost-constrained organizations --- smaller enterprises and
public sector --- should prioritize two things: replacing
unencrypted GRE tunnels with WireGuard (at zero additional
license cost) and migrating remote access to ZTNA, which
scales more predictably than VPN concentrators.

Compliance-driven organizations face a more nuanced situation.
IPSec with IKEv2 and NIST-approved cipher suites remains the
most established option for regulatory frameworks that specify
allowed algorithms. However, those frameworks are beginning
to address PQC migration requirements, and early adoption of
RFC~9370 hybrid key exchange positions an organization ahead
of the curve rather than in a reactive posture.

% ---------------------------------------------------------------
\section{Future Work}

Three directions seem most consequential for this space.

The most pressing is a standardized post-quantum extension for
WireGuard. The Noise handshake protocol that WireGuard uses
does not currently have a standardized PQC variant analogous
to RFC~9370 for IKEv2. Hu's 2020 prototype~\cite{hu2020pq}
demonstrated that hybrid PQC integration with WireGuard is
feasible with moderate overhead, but a standardized extension
is needed before enterprise deployments can rely on it.

A second important direction is formal verification of ZTNA
policy engines. The security of zero trust architectures
depends entirely on the correctness of the policy language
and its enforcement. Misconfiguration is a real risk that is
difficult to detect without formal methods --- a policy that
appears correct can silently grant more access than intended.
Model-checking approaches applied to ZTA policy languages
would provide assurance that current tooling cannot.

Third, post-quantum zero-knowledge proofs for authentication
represent an important gap at the intersection of privacy and
quantum resistance. Current ZKP schemes used in privacy-preserving
authentication depend on elliptic curve pairings that are
vulnerable to Shor's algorithm. Hash-based alternatives like
STARKs and lattice-based ZKPs offer post-quantum properties
but need further development and standardization before they
can replace current schemes at enterprise scale.

% ---------------------------------------------------------------
\section{Conclusion}
\label{sec:conclusion}

Enterprise VPN security is not failing because the underlying
protocols were badly designed. GRE, IPSec, and MPLS were
designed for the threat environment of the 1990s and early
2000s, and they work as specified. The problem is that the
threat environment has changed in ways those specifications
did not anticipate: adversaries who collect encrypted traffic
for future quantum decryption, and attackers who exploit the
implicit post-authentication trust that every perimeter-based
VPN extends to connected devices.

This paper addressed both gaps in sequence. The comparative
survey in Sections~\ref{sec:relatedwork}--\ref{sec:threatscape}
characterized the current state of six enterprise connectivity
technologies and made explicit the threat dimensions most
relevant to 2024--2026: lateral movement risk, quantum key
vulnerability, and performance overhead. The QRA-ZTA framework
in Section~\ref{sec:framework} proposed a response: hybrid
PQC key exchange per NIST FIPS-203 and RFC~9370 at the
cryptographic layer, eBPF/XDP enforcement in the data plane,
NIST SP~800-207 zero trust verification in the trust layer,
and AI-driven behavioral anomaly detection in the intelligence
layer. Each stratum targets a distinct gap, and together they
can be layered incrementally onto existing enterprise
infrastructure.

The core argument is that organizations do not have to choose
between protecting against quantum threats and protecting
against lateral movement. QRA-ZTA shows that both can be
addressed together, using mechanisms that are either
standardized today or close to standardization, in a way
that respects the operational constraints of enterprises that
have built significant infrastructure around the technologies
surveyed here.

\bibliographystyle{IEEEtran}
\bibliography{references_final}

\end{document}
